% answers_short.tex —— 2003 年·上海卷·理科 · 简短答案速查
% 由 tools/make_short_answers.py 自 answers.tex 的 \gksolhead/\gkanswer 提取，
% 并与 2003/2003_数学真题_上海理_答案要点.md 逐题比对后生成，勿手改（改源 MD 后重跑）。

\begin{gknotice}
\textbf{使用说明}：本节只给最终答案便于快速核对，完整推导见解析版。答案为编写者独立推导，\textbf{非官方标准答案}。
\end{gknotice}

\begin{gknotice}
选择题答案中的 U 表示待核，未给出猜测答案，详情见解析版。
\end{gknotice}

\gkpart{填空题}
\begin{itemize}[leftmargin=2.6em]
  \item[\textbf{1.}] \(\pi\)
  \item[\textbf{2.}] \(\dfrac{4\pi}{3}\)
  \item[\textbf{3.}] \(-49\)
  \item[\textbf{4.}] \(\left(\dfrac{\sqrt2}{2},\dfrac{3\pi}{4}\right)\)
  \item[\textbf{5.}] \(\arccos\dfrac1{\sqrt5}\)
  \item[\textbf{6.}] \([1,3]\)
  \item[\textbf{7.}] \(\arccos\dfrac{11}{16}\)
  \item[\textbf{8.}] 例如 \((a_1,q)=\left(1,\dfrac12\right)\)
  \item[\textbf{9.}] \(\dfrac{119}{190}\)
  \item[\textbf{10.}] \(2.6\)
  \item[\textbf{11.}] \(4\pi\)
  \item[\textbf{12.}] 学生解答不正确；\(|PF_2|=1\) 不可实现，正确结果为 \(17\)。
\end{itemize}

\gkpart{选择题}
\begin{center}\begin{tabular}{ccccc}
\toprule
题号 & 13 & 14 & 15 & 16 \\
\midrule
答案 & C & U & D & B \\
\bottomrule
\end{tabular}\end{center}

\gkpart{解答题}
\begin{itemize}[leftmargin=2.6em]
  \item[\textbf{17.}] 最小值 \(\sqrt2\)，最大值 \(\dfrac32\)
  \item[\textbf{18.}] 体积 \(V=8\sqrt3\)
  \item[\textbf{19.}] (1) 两个和依次为 \(a_1(1-q)^2\)、\(a_1(1-q)^3\)；(2) \(\sum_{k=0}^{n}(-1)^k\binom nk a_{k+1}=a_1(1-q)^n\)
  \item[\textbf{20.}] (1) \(l=\dfrac{88}{\sqrt7}\approx33.3\text{ m}\)；(2) \(h=4.5\sqrt2\approx6.4\text{ m},\quad l=22\sqrt2\approx31.1\text{ m}\)
  \item[\textbf{21.}] (1) \(\overrightarrow{AB}=(6,8)\)；(2) \((x-1)^2+(y-3)^2=10\)；(3) 存在，且 \(a>\dfrac32\)
  \item[\textbf{22.}] (1) \(f(x)=x\notin M\)；(2) \(a^x\in M\)；(3) \(k\in\pi\mathbb Z\)
\end{itemize}
